% General academic CV (also used for TA applications).
% Built by `make` into assets/files/CV.pdf, which the website links to.
\documentclass{clean_cv}
\hypersetup{pdftitle={Yu Duan - CV}, pdfauthor={Yu Duan}}

\newcommand{\reviewvenue}{Under review}
% Single source of truth for publications. Each document \input's this file and
% then picks/orders entries with \pubX inside a publist environment.
%
% Each document defines \reviewvenue before \input'ing this file, e.g.
%   \newcommand{\reviewvenue}{Under review}            (public academic CV)
%   resumes read it from ../private/reviewvenue.tex (gitignored) if present, so the
%   venue of double-blind submissions never appears in this public repo.

\newcommand{\me}{\textbf{Yu Duan}}

% ---- Conference papers ------------------------------------------------------
\newcommand{\pubPOCO}{\pub{https://openreview.net/forum?id=Xma2G1ak3H}
  {POCO: Scalable Neural Forecasting through Population Conditioning}
  {\me, Hamza Tahir Chaudhry, Misha B. Ahrens, Christopher D. Harvey, Matthew G. Perich, Karl Deisseroth$^\dagger$, Kanaka Rajan$^\dagger$}
  {Advances in Neural Information Processing Systems (NeurIPS), 2025; Computational and Systems Neuroscience (COSYNE), 2026}\label{pub:poco}}

\newcommand{\pubHebbian}{\pub{https://openreview.net/forum?id=2WklawyeI08}
  {Hebbian and Gradient-based Plasticity Enables Robust Memory and Rapid Learning in RNNs}
  {\me, Zhongfan Jia, Qian Li, Yi Zhong, Kaisheng Ma}
  {International Conference on Learning Representations (ICLR), 2023}\label{pub:hebbian}}

% ---- Under review -----------------------------------------------------------
\newcommand{\pubSlider}{\pub{}
  {SliderField: Fast Dense Sliders via Gaussian Random Fields}
  {\me, Silong Yong, Soo Ye Kim, Yilin Wang, Zhixin Shu, Tianyu Wang, Mang Tik Chiu, Xu Ma, Ayush Shrivastava, Ding Liu, Lingzhi Zhang, Yuqian Zhou, Eli Shechtman, Sohrab Amirghodsi, Xiaohang Zhan}
  {\reviewvenue, 2026}\label{pub:slider}}

\newcommand{\pubVerify}{\pub{}
  {Agentic Verification for Faithful Image Editing}
  {Silong Yong, \me, Yilin Wang, Mang Tik Chiu, Ayush Shrivastava, Zhixin Shu, Yuqian Zhou, Xu Ma, Ding Liu, Lingzhi Zhang, Eli Shechtman, Sohrab Amirghodsi, Yaqi Xie, Katia P. Sycara, Xiaohang Zhan}
  {\reviewvenue, 2026}\label{pub:verify}}

% Short-author variants for 1-page resumes
\newcommand{\pubSliderShort}{\pub{}
  {SliderField: Fast Dense Sliders via Gaussian Random Fields}
  {\me, Silong Yong, Soo Ye Kim, \ldots, Eli Shechtman, Sohrab Amirghodsi, Xiaohang Zhan}
  {\reviewvenue, 2026}}

\newcommand{\pubVerifyShort}{\pub{}
  {Agentic Verification for Faithful Image Editing}
  {Silong Yong, \me, Yilin Wang, \ldots, Katia P. Sycara, Xiaohang Zhan}
  {\reviewvenue, 2026}}

% ---- Abstracts --------------------------------------------------------------
\newcommand{\pubCosyneCausal}{\pub{}
  {Multi-animal Models of Neural Dynamics Predict Responses to Unseen Perturbations}
  {\me, Hamza Tahir Chaudhry, Kayvon Daie, Matthew G. Perich, Kanaka Rajan}
  {Submitted to Computational and Systems Neuroscience (COSYNE), 2027}\label{pub:causal}}

\newcommand{\pubCosyneDist}{\pub{}
  {The Average Never Happens: Distributional Tests for Neural Dynamics Models}
  {\me, Matthew G. Perich, Hamza Tahir Chaudhry, Kanaka Rajan}
  {Submitted to Computational and Systems Neuroscience (COSYNE), 2027}\label{pub:dist}}

\newcommand{\pubNMRNN}{\pub{}
  {Modeling Rapid Neuromodulation in the Cortex-Basal Ganglia-Thalamus Loop}
  {Julia Costacurta$^*$, \me$^*$, Kanaka Rajan, Scott Linderman, John Assad}
  {Computational and Systems Neuroscience (COSYNE), 2025}\label{pub:nmrnn}}

\newcommand{\pubWM}{\pub{https://www.biorxiv.org/content/10.1101/2023.03.30.534982v1}
  {Natural Constraints Explain Working Memory Capacity Limitations in Sensory-Cognitive Models}
  {Yudi Xie$^*$, \me$^*$, Aohua Cheng, Pengcen Jiang, Christopher Cueva, Guangyu Robert Yang}
  {Computational and Systems Neuroscience (COSYNE), 2023; Cognitive Computational Neuroscience (CCN), 2022}\label{pub:wm}}

% ---- Compact one-line versions for 1-page resumes ---------------------------
\newcommand{\cpubPOCO}{\cpub{https://openreview.net/forum?id=Xma2G1ak3H}{POCO: Scalable Neural Forecasting through Population Conditioning}{\textbf{Y.~Duan} et al.}{NeurIPS 2025; COSYNE 2026}}
\newcommand{\cpubHebbian}{\cpub{https://openreview.net/forum?id=2WklawyeI08}{Hebbian and Gradient-based Plasticity Enables Robust Memory and Rapid Learning in RNNs}{\textbf{Y.~Duan} et al.}{ICLR 2023}}
\newcommand{\cpubSlider}{\cpub{}{SliderField: Fast Dense Sliders via Gaussian Random Fields}{\textbf{Y.~Duan} et al.}{\reviewvenue}}
\newcommand{\cpubVerify}{\cpub{}{Agentic Verification for Faithful Image Editing}{S.~Yong, \textbf{Y.~Duan} et al.}{\reviewvenue}}
\newcommand{\cpubCosyneCausal}{\cpub{}{Multi-animal Models of Neural Dynamics Predict Responses to Unseen Perturbations}{\textbf{Y.~Duan} et al.}{COSYNE 2027 (submitted)}}
\newcommand{\cpubCosyneDist}{\cpub{}{The Average Never Happens: Distributional Tests for Neural Dynamics Models}{\textbf{Y.~Duan} et al.}{COSYNE 2027 (submitted)}}
\newcommand{\cpubNMRNN}{\cpub{}{Modeling Rapid Neuromodulation in the Cortex-Basal Ganglia-Thalamus Loop}{J.~Costacurta$^*$, \textbf{Y.~Duan}$^*$ et al.}{COSYNE 2025}}
\newcommand{\cpubWM}{\cpub{https://www.biorxiv.org/content/10.1101/2023.03.30.534982v1}{Natural Constraints Explain Working Memory Capacity Limitations in Sensory-Cognitive Models}{Y.~Xie$^*$, \textbf{Y.~Duan}$^*$ et al.}{COSYNE 2023}}

% Shared bits used by all CV / resume variants.

\newcommand{\linkedinurl}{https://www.linkedin.com/in/yu-duan-a15617308/}

% Contact line for resumes (plain-text URLs are easier for ATS parsers and print)
\newcommand{\resumecontact}{%
\begin{center}\small
    \faCenter{envelope}~\href{mailto:duany19@mit.edu}{duany19@mit.edu} \enspace
    \faCenter{globe}~\href{https://yuvenduan.github.io}{yuvenduan.github.io} \enspace
    \faCenter{github}~\href{https://github.com/yuvenduan}{github.com/yuvenduan} \enspace
    \faCenter{linkedin}~\href{\linkedinurl}{linkedin.com/in/yu-duan-a15617308}
\end{center}\vspace{-4pt}}

% Honors in the academic (date-column) layout
\newcommand{\honorslist}{%
\dateentry{2022}{Scholarship for Scientific Innovation, Tsinghua University}
\dateentry{2021}{Scholarship for Academic Excellence, Tsinghua University}
\dateentry{2019}{Freshman Scholarship, Tsinghua University}
\dateentry{2017}{Gold Medal, National Olympiad in Informatics (NOI), China}}

% Honors for resumes
\newcommand{\honorsresume}{%
\begin{bullets}
\item Gold Medal, National Olympiad in Informatics (NOI), China \hfill 2017
\item Scholarship for Scientific Innovation, Tsinghua University \hfill 2022
\item Scholarship for Academic Excellence, Tsinghua University \hfill 2021
\item Freshman Scholarship, Tsinghua University \hfill 2019%
\end{bullets}}

% Skills (academic CV); resumes define their own, targeted skills lists.
% TODO: add tools you actually use (e.g. multi-GPU / distributed training, SLURM, Hugging Face, W&B).
\newcommand{\skillsline}{%
\textbf{Programming:} Python, C/C++; PyTorch, JAX, scikit-learn. \\
\textbf{Methods:} transformers, diffusion and flow matching, post-training (SFT / RL), RNNs, neural data analysis.}


\author{Yu Duan}

\begin{document}

\maketitle
\begin{center}
    \faCenter{envelope}~\href{mailto:duany19@mit.edu}{duany19@mit.edu} \quad
    \faCenter{globe}~\href{https://yuvenduan.github.io}{yuvenduan.github.io} \quad
    \faCenter{github}~\href{https://github.com/yuvenduan}{github.com/yuvenduan} \quad
    \faCenter{linkedin}~\href{\linkedinurl}{LinkedIn}
\end{center}

\section{Education}
\begin{datetabular}{9.6em}
\dateentry{Sep 2023 -- present}{
\textbf{Massachusetts Institute of Technology}, Cambridge, MA \\
Ph.D. in Electrical Engineering and Computer Science \\
Research advisor: Prof.\ Kanaka Rajan (Kempner Institute, Harvard University)
}
\dateentry{Sep 2019 -- Jun 2023}{
\textbf{Tsinghua University}, Beijing, China \\
B.Eng.\ in Computer Science and Technology \\
Institute for Interdisciplinary Information Sciences
}
\dateentry{Feb 2022 -- Jul 2022}{
\textbf{Massachusetts Institute of Technology}, Cambridge, MA \\
Visiting Student, Department of Brain and Cognitive Sciences
}
\end{datetabular}

\section{Research Experience}
\begin{datetabular}{9.6em}

\dateentry{Feb 2024 -- present}{
\textbf{Foundation Models of Neural Dynamics} \\
\textit{Advisor: Prof.\ Kanaka Rajan (Harvard); with Prof.\ Karl Deisseroth (Stanford)}
\begin{itemize}[leftmargin=1.1em]
\item Neural forecasting model pretrained across sessions, animals and species \ref{pub:poco}.
\item Multi-animal models that predict responses to unseen optogenetic stimulation \ref{pub:causal}.
\item Flow-matching models and distributional tests for stochastic neural dynamics \ref{pub:dist}.
\end{itemize}\eatvspace}

\dateentry{May 2026 -- Aug 2026}{
\textbf{Post-training and Efficient Sampling for Image Editing} \\
\textit{Research Scientist Intern, Adobe}
\begin{itemize}[leftmargin=1.1em]
\item Supervised fine-tuning (SFT) of a foundation model for continuous physical edits.
\item Training-free 4--5$\times$ acceleration of dense editing sliders in flow models \ref{pub:slider}.
\item Agentic verifier for image editing, used as an RL reward and data filter \ref{pub:verify}.
\end{itemize}\eatvspace}

\dateentry{Feb 2022 -- Dec 2023}{
\textbf{Multi-system Models of Working Memory} \\
\textit{Advisor: Prof.\ Guangyu Robert Yang, Brain and Cognitive Sciences, MIT} \\
Human-like capacity limits from natural constraints in sensory-cognitive models \ref{pub:wm}.
}

\dateentry{Mar 2023 -- Jun 2023}{
\textbf{Predicting fMRI Responses with Vision-Language Models} \\
\textit{Undergraduate thesis. Advisor: Prof.\ Pei Sun, Department of Psychology, Tsinghua University}
}

\dateentry{Mar 2021 -- Sep 2022}{
\textbf{Plastic Recurrent Networks and Models of the Fly Mushroom Body} \\
\textit{Advisors: Prof.\ Kaisheng Ma and Prof.\ Yi Zhong, Tsinghua University} \\
Learned Hebbian plasticity in RNNs for robust memory and rapid learning \ref{pub:hebbian}.
}

\end{datetabular}

\section{Publications}
\subsection{Conference Papers}\hfill\par
\begin{publist}
\pubPOCO
\pubHebbian
\end{publist}

\subsection{Manuscripts Under Review}\hfill\par
\begin{publist}[resume]
\pubSlider
\pubVerify
\end{publist}

\subsection{Conference Abstracts}\hfill\par
\begin{publist}[resume]
\pubCosyneCausal
\pubCosyneDist
\pubNMRNN
\pubWM
\end{publist}
\hfill{\small $^*$Equal contribution. $^\dagger$Co-senior authors.}

% TODO (TA applications): add teaching / mentoring experience here if you have any, e.g.
% \section{Teaching and Mentoring}
% \begin{datetabular}{9.6em}
% \dateentry{Fall 2025}{\textbf{Teaching Assistant}, 6.XXXX Course Name, MIT \\ Led recitations for N students; held office hours; wrote problem sets}
% \end{datetabular}

\section{Honors and Awards}
\begin{datetabular}{9.6em}\setlength{\itemsep}{0pt}
\honorslist
\end{datetabular}

\section{Skills}
\skillsline

\section{Service}
\textbf{Reviewer:} International Conference on Learning Representations (ICLR)

\end{document}
